% Author: SuFan; MIT License. Illustrative energy-use comparison, not measured data.
\documentclass{article}
\usepackage[margin=20mm]{geometry}
\usepackage{gradbars}
\pagestyle{empty}
\begin{document}
% BEGIN TABLE
\begingroup
\begin{center}
\begin{gradtable}[width=\linewidth,align=c,header align=c,
  column sep=3pt,row stretch=1.7,label width=9mm]
\gradcolumn[width=22mm]{building}{Building}
\gradcolumn[type=number,width=16mm,
  options={precision=0}]{before}{Before}
\gradcolumn[type=number,width=16mm,
  options={precision=0}]{after}{After}
\gradcolumn[type=dumbbell,options={min=0,max=200,
  precision=0,compare label=delta,track color=black!35,
  stem width=.9pt}]{change}{Change}
\gradheader{\gradspan{1}{Site}
  \gradspan{2}{Annual use (kWh/m$^2$)}
  \gradspan{1}{Common scale: 0--200}}
\gradrow{Atrium,160,120,160/120}
\gradrow{Studio,130,90,130/90}
\gradrow{Archive,80,76,80/76}
\gradrow{Annex,110,128,110/128}
\end{gradtable}
\end{center}
\endgroup
% END TABLE
\noindent Hollow points are before, filled points are after. All dumbbells
share 0--200. Delta is after minus before, in kWh/m$^2$; a negative value means
lower annual use. These are illustrative comparisons, not causal estimates.
\end{document}
